\section{Further research questions and explicit conjectures}
\label{sec:research-agenda}

The results isolate several precise next problems. The first group
would complete finite-parameter statements already suggested by the
source reports. The later groups ask for uniform transitions or new
structures beyond the results proved here.

\subsection{Complete the finite-index Euler extremum problem}

\begin{conjecture}[Axis reduction at every truncation]
\label{conj:all-index-axis}
For every integer $N\geq1$,
\[
 C_N=\max_{b>0}R_N(0,b),
\]
and the unique maximizer on $a\geq0$, $b>0$ lies on that axis.
\end{conjecture}

The case $N=1$ is part of the known universal-constant theorem;
the present article proves all sufficiently large $N$. The conjecture
is about a global maximum over both orders. It should not be replaced
by a pointwise comparison $R_N(a,b)\leq R_N(0,b)$ at every fixed $b$:
that stronger statement is not the claim and has no such proof here.
A productive approach would combine exact small-index analysis with
an effective, moderately sized threshold in the localization argument.
The relevant difficulties occur outside the eventual competitive band.

An explicit threshold $N_0$ is itself useful. The present split at
$b=2\log N$ produces a slowly vanishing factor
$N^{1-2\log(5/3)}$. Improving this estimate, perhaps with a more
carefully optimized gamma lower-tail bound, would turn the eventual
theorem into a practical enclosure theorem. A finite numerical scan
without a proved tail threshold would not close this question.

\begin{conjecture}[All-index kernel uniqueness and decrease]
\label{conj:all-index-kernel}
For every $N\geq2$, $K_N$ has a unique global maximizer in $(0,1)^2$
and $M_{N+1}<M_N$.
\end{conjecture}

This restates the remaining portion of the sharp-Euler report's
conjecture. The exact $N=2$ elimination and the analytic branch near
$1/N=0$ suggest combining algebraic continuation with certified
boundary exclusions. Monotonicity of the envelope $B_N$ alone does
not prove monotonicity of $M_N$.

\subsection{Refine the gamma-transition asymptotics}

The next Euler scale has the form
$N^{-\mu}(\log N)^{-1/2}$ and carries a large first inverse-logarithmic
coefficient. Determine a uniformly useful approximation across the
moderate-$N$ regime and the true asymptotic regime. The exact positive
remainder integral is a better starting point than subtraction of
nearly equal Euler sums. One target is a resummed transition integral
whose error remains explicit while retaining the slowly decaying tail
responsible for the large coefficient.

For other residual kernels, moment competition may put the analogous
Mellin exponent at an integer. Classify when logarithmic resonances
replace a pure fractional-power correction, and determine the resulting
maximizer displacement. The present proof suggests a general method:
identify the competing exponential moments, localize the parameters,
and analyze the positive remainder near the moving kernel transition.
The exponent and amplitude need to be derived anew for each kernel.

The joint profile in \cref{eul:thm:joint-profile} gives the exact first
loss from leaving the axis. Higher joint terms could quantify stability
of near-optimal order selection at prescribed precision. They should
be stated on explicit joint scales for $a$, $b-\widehat b_N$, and $N$;
pointwise small-$a$ and large-$N$ expansions are not automatically uniform.

\subsection{Determine the arithmetic of the oscillation phase}

For a fixed algebraic value such as $A=2$, is
$\theta(A)/(2\pi)$ irrational? The article proves sign-density
statements in both the rational and irrational cases, so no unproved
phase arithmetic is needed for those theorems. A proof of irrationality
at $A=2$ would give density $1/2$ for each sign and would make rotation
methods available for more precise gap statistics.

Find the least eventual block length that must contain both signs.
The bound $\lfloor\pi/\theta\rfloor+2$ is explicit and gives four at
$A=2$. Whether a smaller block can work is a question about the actual
rotation orbit and its distances from the sign boundaries, together
with the coefficient error. Likewise, quantify the number of indices
where the normalized velocity is unusually small. A relative cosine
approximation is unsuitable for these exceptional indices; the
all-order additive expansion is the available tool.

The finite-part identity for the critically weighted first derivative
can be extended to higher derivative sums. Determine their full
subtraction polynomials in fractional powers of the cutoff, and
identify which finite constants are local Puiseux coefficients.
This should give explicit renormalized Stirling identities, with
ordinary convergence and regularized limits carefully distinguished.

\subsection{A proposed bulk quantization formula}

Let $L_{n,j}$ be the positive zeros of $p_n$ in increasing order,
$\theta_{n,j}=\theta(e^{L_{n,j}})$, and evaluate $u_*,\chi,\chi_1$
at $A=e^{L_{n,j}}$. Put $\phi=-\arg\chi$ as before.
The following indexed refinement is consistent with the finite root
isolations supplied in the package.

\begin{conjecture}[An indexed phase formula in the bulk]
\label{conj:bulk-quantization}
For every fixed $\epsilon\in(0,1/2)$, uniformly for
$\epsilon(n-1)\leq j\leq(1-\epsilon)(n-1)$,
\begin{equation}\label{eq:bulk-quantization-conjecture}
 n\theta_{n,j}+\phi(\theta_{n,j})
 =\pi\left(j+\frac12\right)
   +\frac1n\Im\left(\frac{\chi_1}{\chi}\right)
   +O_\epsilon(n^{-2}).
\end{equation}
\end{conjecture}

The unproved content includes the global root indexing, in particular
the exact shift $j+1/2$. Local transfer locates roots near a half-integer
phase and gives the displayed correction once the index is identified;
the bulk counting theorem alone does not fix that index. One possible
route is to match the Bessel endpoint asymptotic to a uniform interior
expansion in an overlap where $j\to\infty$ and $j/n\to0$.
This is a conjectural asymptotic identity, not a formula used anywhere
in the proofs of this article.

\subsection{Connect the smallest, bulk, and largest root regimes}

The three proved regimes are
\[
 L_{n,j}\asymp n^{-2}\quad(j\text{ fixed from the small end}),
 \qquad L_{n,j}=O(1)\quad(j/n\text{ in the bulk}),
\]
and $X_{n,k,j}\sim n/j$ for $j$ fixed from the large end.
Develop uniform expansions when the root index grows while remaining
$o(n)$. At the large end, the main question is how the expansion around
the $j$th zero of $1/\Gamma(1-w)$ changes when $j$ is no longer fixed.
At the small end, the Bessel zero expansion should be matched to the
singularity-transfer phase. Either transition would sharpen the weak
zero law and give useful discrepancy estimates.

The all-orders largest-root theorem treats fixed positive integer $k$.
Determine uniform ranges of $k=k(n)$, and carefully extend the theorem
to Cohen's derivative interpretation for nonpositive $k$ if desired.
The exact logarithmic change of variables remains a promising starting
point, but the compact-parameter gamma expansion used here does not
justify growing-$k$ or growing-$j$ conclusions.

Recent polynomial-profile methods \cite{JalowyKabluchkoMarynych2026}
may supply another description of the bulk law. Determine whether the
Lambert polynomial family admits a useful limiting Cauchy transform,
free-convolution interpretation, or a potential-theoretic variational
principle. Any such identification should explicitly handle its
degree-dependent differential operator; the fixed-function examples
in the general theory are not automatic substitutes.

\subsection{Return the analytic tools to polylogarithmic identities}

The exact error kernels remain available as certified numerical tools
for cyclotomic and Gaussian polylogarithm identities. The unresolved
mixed $S_6$ reduction is an appropriate arithmetic target, but a
candidate relation requires an independent reduction mechanism such
as iterated-integral relations, a differential identity with controlled
boundary values, or a valid coaction argument. A high-precision integer
relation search alone would still leave a conjecture.

For roots of unity beyond $i$, seek a sector or matrix positivity
statement that replaces the scalar order of the present Gaussian
kernel. At higher depth, identify the correct iterated compensated
differences and prove integrability at each step. The all-depth
oscillation results already in the incoming material concern a
different part of this problem; they should be used with their stated
hypotheses, not treated as a ready-made extension of the present
two-order Euler extremum theorem.

Finally, the exact finite components provide natural formalization
targets: the Euler coefficient identity, the rational derivative and
resultant certificates, finite-field irreducibility, the elementary
real-rooted operators, and the algebra of the first edge coefficients.
A later formal development could then connect them to gamma
integration, Rouch\'e continuation, and singularity transfer. The
package supplies reproducible ordinary proofs and finite arithmetic;
that proposed formal bridge remains future work.
